\section*{Bloque C. Casos con implicancias regulatorias y tarifarias}

\noindent\textit{Preguntas 66 a 90. Aplican el análisis del monopolio y del monopolio natural a los servicios de agua y saneamiento, la regulación tarifaria, la discriminación de precios y la pérdida de eficiencia.}

\q{Una empresa de agua que abastece toda una ciudad y presenta costos medios decrecientes en el rango relevante de demanda es:}
\begin{opts}
    \item un monopolio natural.
    \item una empresa en competencia perfecta.
    \item un monopsonio.
    \item una franja competitiva.
\end{opts}

\q{Si el regulador fija $P=CMg$ en un monopolio natural, la empresa:}
\begin{opts}
    \item obtiene beneficios extraordinarios positivos permanentes.
    \item obtiene un beneficio económico nulo en el largo plazo.
    \item sufre pérdidas que requieren subsidio o cobertura.
    \item abandona de inmediato la operación por falta de renta.
\end{opts}

\q{La pérdida de eficiencia de un monopolio no regulado se debe a que:}
\begin{opts}
    \item fija el precio igual al costo marginal.
    \item produce más que el óptimo socialmente eficiente.
    \item carece de costos fijos y marginales positivos.
    \item fija $P>CMg$ produciendo bajo el óptimo social.
\end{opts}

\q{Un precio máximo igual al costo medio permite:}
\begin{opts}
    \item generar beneficios extraordinarios permanentes.
    \item arrojar pérdidas operativas inevitables cada año.
    \item alcanzar beneficio económico nulo cubriendo costos.
    \item impedir la recuperación contable del capital fijo.
\end{opts}

\q{Una tarifa en dos partes en el servicio de agua combina:}
\begin{opts}
    \item exclusivamente un cargo fijo mensual sin medición.
    \item un cargo fijo más un cargo por unidad consumida.
    \item únicamente un cargo volumétrico por metro cúbico.
    \item un impuesto pigouviano uniforme sobre el consumo.
\end{opts}

\q{La discriminación de precios de tercer grado en agua puede aplicarse entre:}
\begin{opts}
    \item todos los usuarios urbanos pagando idéntica tarifa.
    \item grupos distintos: residencial, comercial e industrial.
    \item un único cliente conectado mediante tarifa plana.
    \item consumidores sin segmentación ni elasticidad distinta.
\end{opts}

\q{Para que la discriminación de precios sea viable se requiere:}
\begin{opts}
    \item estructura de competencia perfecta con muchos ofertantes.
    \item libre entrada y salida instantánea de competidores.
    \item venta de un producto homogéneo con libre arbitraje.
    \item poder de mercado, segmentación y ausencia de arbitraje.
\end{opts}

\q{Si la demanda residencial es más inelástica que la comercial, bajo discriminación de tercer grado la tarifa residencial será:}
\begin{opts}
    \item menor.
    \item igual.
    \item mayor.
    \item cero.
\end{opts}

\q{La discriminación de precios de primer grado en agua:}
\begin{opts}
    \item maximiza el excedente neto retenido por el consumidor.
    \item carece por completo de impactos distributivos relevantes.
    \item genera un beneficio económico nulo para el monopolista.
    \item captura todo excedente del usuario, afectando equidad.
\end{opts}

\q{Un monopolio natural se define por:}
\begin{opts}
    \item costos medios decrecientes en el rango de demanda.
    \item costos medios crecientes desde niveles iniciales.
    \item competencia perfecta con libertad de entrada plena.
    \item una curva de oferta agregada de largo plazo horizontal.
\end{opts}

\q{La regulación por tasa de retorno puede incentivar:}
\begin{opts}
    \item sobreinversión o sobrecapitalización.
    \item desinversión total en infraestructura.
    \item reducción de costos marginales a cero.
    \item competencia perfecta entre operadores.
\end{opts}

\q{Un precio máximo o price cap incentiva:}
\begin{opts}
    \item el aumento deliberado de costos operativos.
    \item la ausencia total de incentivos de gestión.
    \item la reducción de costos y mayor eficiencia.
    \item la salida inmediata del operador del mercado.
\end{opts}

\q{La competencia por el mercado, mediante una subasta de concesión, puede:}
\begin{opts}
    \item suprimir por completo cualquier marco regulatorio.
    \item elevar artificialmente los costos operativos medios.
    \item consolidar un monopolio privado sin supervisión.
    \item disciplinar al operador compitiendo por el mercado.
\end{opts}

\q{La competencia por comparación o yardstick competition consiste en:}
\begin{opts}
    \item coordinar tarifas uniformes entre operadores rivales.
    \item estatizar íntegramente los activos de saneamiento.
    \item pactar precios regionales mediante acuerdos privados.
    \item comparar empresas similares para fijar metas o tarifas.
\end{opts}

\q{La producción pública o municipal del servicio de agua:}
\begin{opts}
    \item resulta invariablemente eficiente sin riesgos de gestión.
    \item es una opción válida con riesgos como el presupuesto blando.
    \item carece por completo de costos de capital y depreciación.
    \item garantiza resultados idénticos a la competencia perfecta.
\end{opts}

\q{La pérdida de eficiencia del monopolio se puede reducir mediante:}
\begin{opts}
    \item mayores barreras legales a la entrada.
    \item regulación, competencia o subastas de concesión.
    \item discriminación perfecta de precios sin control.
    \item reducción forzosa de la demanda residencial.
\end{opts}

\q{Un subsidio al consumo de agua para hogares vulnerables se justifica principalmente por:}
\begin{opts}
    \item equidad distributiva, aunque pueda introducir distorsiones.
    \item eficiencia asignativa pura bajo costo marginal constante.
    \item ausencia de costos fijos en redes de distribución urbana.
    \item eliminación directa de la demanda de usuarios no pobres.
\end{opts}

\q{La tarifa en bloques crecientes es una forma de discriminación de precios:}
\begin{opts}
    \item de primer grado.
    \item de segundo grado.
    \item de tercer grado.
    \item uniforme.
\end{opts}

\q{La tarifa en bloques crecientes puede:}
\begin{opts}
    \item carece de efecto alguno sobre consumo o recaudación.
    \item resulta invariablemente eficiente en primer mejor.
    \item suprime toda pérdida irrecuperable de bienestar social.
    \item balancear conservación, equidad de pago y eficiencia.
\end{opts}

\q{El cargo fijo en la tarifa de agua permite:}
\begin{opts}
    \item cubrir costos fijos reduciendo distorsión del cargo variable.
    \item elevar artificialmente el costo marginal de distribución.
    \item suprimir el consumo básico de los usuarios residenciales.
    \item fijar obligatoriamente el precio variable unitario en cero.
\end{opts}

\q{La regulación del monopolio natural enfrenta un trade-off entre:}
\begin{opts}
    \item inflación monetaria y tasa de desempleo.
    \item equidad vertical y equidad horizontal.
    \item eficiencia asignativa y sostenibilidad financiera.
    \item precio promedio histórico y demanda fija.
\end{opts}

\q{El monopolio no regulado en agua tiende a:}
\begin{opts}
    \item fijar $P>CMg$ produciendo bajo el óptimo.
    \item fijar $P=CMg$ cubriendo sus costos totales.
    \item producir en exceso respecto al óptimo social.
    \item establecer una tarifa volumétrica igual a cero.
\end{opts}

\q{La solución de primer mejor en un monopolio natural requiere:}
\begin{opts}
    \item $P=CMg$ con subsidio para cubrir pérdidas.
    \item $P=CMe$ sin transferencias del Tesoro.
    \item $P>CMe$ generando rentas extraordinarias.
    \item ausencia completa de regulación tarifaria.
\end{opts}

\q{La solución de segundo mejor en un monopolio natural suele ser:}
\begin{opts}
    \item $P=CMg$ cubierto con subsidios directos.
    \item $P=0$ financiado con deuda pública.
    \item $P=CMe$ o precios óptimos de Ramsey.
    \item el precio libre de monopolio no regulado.
\end{opts}

\q{La discriminación de precios en agua puede:}
\begin{opts}
    \item resultar siempre eficiente sin efectos redistributivos.
    \item carecer de todo impacto sobre cobertura y facturación.
    \item mejorar recuperación de costos, aunque afecte equidad.
    \item eliminar la necesidad de supervisión regulatoria estatal.
\end{opts}
